\vspace{-10pt}
\section{静磁场中的磁偶极子}
我们知道，对于匀强磁场中的磁偶极子，也就是载流线圈，根据\Refx{定律}{匀强磁场中的安培力性质}可知其受到的磁场力为零，但是磁偶极子在匀强磁场中仍会受到磁力矩的作用。
\begin{Figure}{静磁场中的磁偶极子}
    \inputfigure[scale=0.6]{Chapter10F_Figure01}\vspace{-10pt}
\end{Figure}\vspace{-5pt}
为了研究的方便，我们先假定这里的磁偶极子是一个矩形载流线圈，且当磁偶极矩的方向确定时，该矩形线圈围需绕磁偶极矩旋转至能使得矩形线圈的两条边垂直于磁场的角度，在这种情况下，两条边上的电流垂直纸面（叉点），两条边上的电流平行纸面（灰色实虚线）。

对于平行纸面的两边上的电流，容易看出两者所受到安培力分别垂直纸面向外和向内，大小相等，方向相反，且通过矩形线圈的中心，既不会产生合力，也不会产生合力矩。

对于垂直直面的两边上的电流，虽然两者的安培力是等大反向的，但是两者的安培力并不是共线的，因此会产生力矩的作用，考虑到两者受力的对称性，我们在计算合力矩时只要考虑其中一条边的电流的力矩并乘以二就可以了，这里我们考虑垂直纸面向上的那根电流。

设垂直纸面的电流的导线长为$l$，设垂直纸面的电流的间距为$d$。

根据\Refx{公式}{匀强磁场中的载流直导线}，垂直纸面向上的电流受到的安培力为
\begin{equation*}
    \vb*{F}=I\vb*{l}\times\vb*{B}
\end{equation*}

这里的$\vb*{l}$的方向是垂直直面向上的电流的方向，设其相对于线圈中心的位矢为$\vb*{d}/2$。这里的矩形线圈受到的力矩，是垂直纸面向上的电流产生的力矩的两倍，根据\Refx{定义}{力矩}
\begin{Equation}[磁偶极子的磁力矩1]
    \vb*{M}=2\qty[\frac{\vb*{d}}{2}\times\qty(I\vb*{l}\times\vb*{B})]=\vb*{d}\times\qty(I\vb*{l}\times\vb*{B})
\end{Equation}
根据矢量三重积的结论，通常情况下矢量叉乘是不满足结合律
\begin{align*}
    \vb*{A}\times(\vb*{B}\times\vb*{C})&=(\vb*{A}\cdot\vb*{C})\vb*{B}-(\vb*{A}\cdot\vb*{B})\vb*{C}\\[4mm]
    (\vb*{A}\times\vb*{B})\times\vb*{C}&=(\vb*{C}\cdot\vb*{A})\vb*{B}-(\vb*{B}\cdot\vb*{C})\vb*{A}
\end{align*}

但如果有$\vb*{A}\cdot\vb*{B}=0$且$\vb*{B}\cdot\vb*{C}=0$，上两式的第二项均为零，此时矢量叉乘满足结合律
\begin{equation*}
    \vb*{A}\times(\vb*{B}\times\vb*{C})=
    (\vb*{A}\times\vb*{B})\times\vb*{C}\qquad(\vb*{A}\perp\vb*{B}~~\vb*{B}\perp\vb*{C})
\end{equation*}

而这里的$\vb*{d},\vb*{l}$和$\vb*{l},\vb*{B}$是垂直的，因此式\ref{式:磁偶极子的磁力矩1}可以化为
\begin{Equation}[磁偶极子的磁力矩2]
    \vb*{M}=(I\vb*{d}\times\vb*{l})\times\vb*{B}=(IS\vb*{e}_n)\times\vb*{B}
\end{Equation}
上式中的$S$为线圈面积$S=dl$，而$\vb*{e}_n$恰好是磁偶极矩方向的单位矢量。

根据\Refx{定义}{磁偶极矩}
\begin{BoxEquation}[磁偶极子的磁力矩]
    \vb*{M}=\vb*{m}\times\vb*{B}
\end{BoxEquation}
而对于任意形状的平面线圈构成的磁偶极子，例如我们最为熟悉的圆线圈，我们可以假象其为无数窄矩形线圈的组合\footnote{由于面积上的相等，平面线圈的磁偶极矩是这些载有相同电流的窄矩形线圈的磁偶极矩的和。}，在两个窄矩形线圈的公共边上会有相反流向的电流，其磁场效应必然相会抵消，就像这些公共边不存在一样，而实际上这些公共边确实是不存在的，因此这种设想并不破坏原有的物理图像，因此\Refx{公式}{磁偶极子的磁力矩}总是适用的。

\Refx{公式}{磁偶极子的磁力矩}指出，磁偶极子在匀强磁场中所受\uwave{磁力矩}的方向，是右手四指由磁偶极矩转向磁感强度时右手拇指的方向，而其旋转方向即这里右手四指的转向。

\Refx{公式}{磁偶极子的磁力矩}中，如果我们记$\vb*{m},\vb*{B}$的夹角为$\theta$，那么当$\theta=\pi/2$时磁力矩最大，而当$\theta=0$或$\theta=\pi$时磁力矩为零，但前者是稳定平衡，而后者是不稳定平衡。
\begin{Figure}{磁偶极子的平衡}
    \subfigure[稳定平衡]
    {
        \inputfigure[scale=0.5]{Chapter10F_Figure02}
    }\qquad
    \subfigure[不稳定平衡]
    {
        \inputfigure[scale=0.5]{Chapter10F_Figure03}
    }
\end{Figure}
因此，磁偶极子在磁力矩的作用下，总是会偏转至磁偶极矩与磁感强度一致的方向。
